\chapter{PENDAHULUAN}
\label{bab:pendahuluan}

\section{Latar Belakang}

Desain struktur merupakan salah satu bagian dari keseluruhan proses
perencanaan bangunan. Proses desain dapat dibedakan menjadi dua tahap, tahap
pertama yaitu desain umum yang merupakan peninjauan secara garis besar
keputusan-keputusan desain, misalnya tata letak struktur, geometri atau
bentuk bangunan, dan material, sedangkan tahap kedua adalah desain secara
terinci yang antara lain meninjau tentang penentuan besar penampang balok
dan kolom, luas dan penempatan tulangan yang dibutuhkan pada struktur beton
bertulang dan sebagainya, untuk itu perlu dimodelkan struktur yang akan
didesain sehingga dapat diperoleh besaran gaya dan informasi perpindahan
yang diperlukan dalam menentukan luasan beton beserta penulangannya.

Untuk memodelkan suatu struktur bangunan yang baik dapat dimodelkan sebagai
portal ruang. Portal ruang merupakan suatu struktur yang pada tiap titiknya
memiliki enam perpindahan yang mungkin terjadi, sehingga bukan lagi dua atau
tiga seperti pada struktur rangka atau portal bidang.

Dalam menentukan dimensi balok dan kolom pada struktur beton bertulang harus
memperhatikan masalah kekuatan dan biaya, terutama pada saat ini dimana
harga material semakin mahal, banyak orang yang menunda bahkan membatalkan
proyeknya karena masalah biaya bahan bangunan yang mahal, oleh karena itu
dalam merencanakan struktur bangunan sangat perlu diperhatikan masalah
tersebut. Kekuatan yang dibutuhkan oleh suatu struktur beton bertulang dapat
dicapai dengan memberikan luasan penampang beton dan tulangan yang cukup.
Analisis perhitungan yang dilakukan diharapkan dapat memperoleh hasil yang
aman dan ekonomis. Untuk mendapatkan hasil yang paling murah (optimal) dapat
dicapai dengan menggunakan proses optimasi. Hasil yang didapat merupakan
hasil yang mempunyai harga struktur yang paling murah tetapi tetap masih
mampu mendukung beban struktur dengan aman.

Metoda optimasi semakin berkembang seiring dengan perkembangan komputer.
Dalam bidang teknik sipil, formulasi masalah optimasinya berupa formula
tidak linear, oleh karena itu metoda optimasi yang berkembang dalam bidang
teknik adalah metoda optimasi non-linear.

Dalam menyelesaikan masalah optimasi non-linear, dapat diselesaikan dengan
bantuan kalkulus yang memerlukan informasi turunan, hal ini dapat dilakukan
apabila masalah yang ada belum begitu komplek, untuk masalah yang sangat
komplek dikembangkan metoda yang idenya berasal dari alam, seperti metoda
algoritma genetik yang berasal dari teori Darwin ``Survival to the
fittest'', metoda \textit{simulated annealing} yang didasarkan pada proses
mengerasnya logam, metoda \textit{flexible polyhedron} yang dalam prosesnya
variabel desain akan membentuk suatu segi banyak yang makin lama semakin
mengecil. Dari banyak metoda di atas, sampai saat ini belum ada metoda yang
dapat digunakan untuk menyelesaikan segala kasus dengan memuaskan, karena
masing-masing metoda memiliki keunggulan dan kerugiannya sendiri-sendiri.

Oleh karena itu dalam tugas akhir ini, untuk melakukan optimasi beton
bertulang pada struktur portal ruang digunakan metoda polihedron fleksibel.

\section{Materi Tugas Akhir}

\subsection{Rumusan Masalah}

Proses optimasi adalah proses mencari nilai variabel desain yang memberikan
nilai optimal yang berupa nilai maksimum ataupun nilai minimum dari fungsi
sasaran tanpa melanggar kendala-kendala yang ada.

Secara umum masalah optimasi dapat dirumuskan sebagai:
\begin{equation}
\text{minimumkan } f(x),\quad x \in X
\label{eq:1-1}
\end{equation}
yang memenuhi kendala:
\begin{align}
g_j(x) &\leq 0, \quad j = 1,2,\ldots,n \label{eq:1-2}\\
h_j(x) &= 0, \quad j = 1,2,\ldots,m \label{eq:1-3}
\end{align}
pada persamaan di atas $x$ menyatakan variabel desain, $X$ menyatakan ruang
variabel desain, $f(x)$ adalah fungsi sasaran, $g(x)$ fungsi kendala
pertaksamaan, $h(x)$ fungsi kendala persamaan, $n$ jumlah kendala
pertaksamaan, $m$ jumlah kendala persamaan.

Harga struktur yang minimum merupakan fungsi sasaran dalam optimasi ini,
dengan variabel disainnya terdiri dari volume beton dan berat tulangan,
dalam hal ini tulangan terdiri dari tulangan tarik dan desak untuk balok,
tulangan geser untuk balok, tulangan kolom, tulangan geser untuk kolom.
Kendala pertaksamaan yang harus dipenuhi berupa momen yang terjadi pada
struktur harus lebih kecil atau sama dengan momen yang dapat didukung oleh
masing-masing batang balok maupun kolom, gaya geser yang terjadi pada
struktur harus lebih kecil atau sama dengan gaya geser yang didukung oleh
beton dan tulangan geser pada masing-masing batang balok maupun kolom, gaya
tekan maupun tarik yang terjadi harus lebih kecil atau sama dengan gaya
tekan atau tarik yang didukung oleh kolom, serta lendutan yang terjadi harus
lebih kecil atau sama dengan lendutan ijin.

Optimasi beton bertulang pada struktur portal ruang ini dapat dirumuskan
sebagai:

\medskip\noindent
\textbf{untuk balok}, minimumkan:
\begingroup\footnotesize
\begin{equation}
\begin{aligned}
f_1(&\mathit{lebar}, \mathit{tinggi}, \mathit{diatltrtm}, \mathit{ntltrtm}, \mathit{diatldstm}, \mathit{ntldstm},\\
    &\mathit{diatltrlp}, \mathit{ntltrlp}, \mathit{diatldslp}, \mathit{ntldslp}, \mathit{diatlgs}, \mathit{jaraktlgs}) =\\[2pt]
  &(\mathit{lebar} \times \mathit{tinggi} \times \mathit{panjang\ balok} \times \mathrm{Rp}(\text{harga beton})/\mathrm{m}^3)\\
 +&\left(\tfrac14 \pi\, \mathit{diatltrtm}^2 \times \mathit{ntltrtm} \times 2 \times 0.25 \times \mathit{panjang\ balok} \times (\text{berat tulangan})_{\mathrm{kg/m}^3} \times \mathrm{Rp}(\text{harga tulangan})/\mathrm{kg}\right)\\
 +&\left(\tfrac14 \pi\, \mathit{diatldstm}^2 \times \mathit{ntldstm} \times 2 \times 0.25 \times \mathit{panjang\ balok} \times (\text{berat tulangan})_{\mathrm{kg/m}^3} \times \mathrm{Rp}(\text{harga tulangan})/\mathrm{kg}\right)\\
 +&\left(\tfrac14 \pi\, \mathit{diatltrlp}^2 \times \mathit{ntltrlp} \times 0.5 \times \mathit{panjang\ balok} \times (\text{berat tulangan})_{\mathrm{kg/m}^3} \times \mathrm{Rp}(\text{harga tulangan})/\mathrm{kg}\right)\\
 +&\left(\tfrac14 \pi\, \mathit{diatldslp}^2 \times \mathit{ntldslp} \times 0.5 \times \mathit{panjang\ balok} \times (\text{berat tulangan})_{\mathrm{kg/m}^3} \times \mathrm{Rp}(\text{harga tulangan})/\mathrm{kg}\right)\\
 +&\Big(\tfrac14 \pi\, \mathit{diatlgs}^2 \times \big(2(\mathit{lebar} - 2\,\mathit{selimut\ beton}) + 2(\mathit{tinggi} - 2\,\mathit{selimut\ beton})\big)\\
   &\times (\mathit{panjang\ balok}/\mathit{jaraktlgs}) \times (\text{berat tulangan})_{\mathrm{kg/m}^3} \times \mathrm{Rp}(\text{harga tulangan})/\mathrm{kg}\Big)
\end{aligned}
\label{eq:1-4}
\end{equation}
\endgroup
yang memenuhi kendala pertaksamaan:
\begin{align}
M_{\text{intern}} &\geq M_{\text{extern}} \label{eq:1-5}\\
G_{\text{intern}} &\geq G_{\text{extern}} \label{eq:1-6}\\
\mathit{Lendutan} &\leq \mathit{Lendutan}_{\text{ijin}} \label{eq:1-7}
\end{align}
pada persamaan di atas $f_1$ menyatakan harga balok pada struktur, $lebar$
dan $tinggi$ menyatakan ukuran penampang balok. Variabel desainnya terdiri
atas $diatltrtm$ menyatakan diameter tulangan tarik pada daerah tumpuan,
$ntltrtm$ menyatakan jumlah tulangan tarik pada daerah tumpuan, $diatldstm$
menyatakan diameter tulangan desak pada daerah tumpuan, $ntldstm$
menyatakan jumlah tulangan desak pada daerah tumpuan, $diatltrlp$
menyatakan diameter tulangan tarik pada daerah lapangan, $ntltrlp$
menyatakan jumlah tulangan tarik pada daerah lapangan, $diatldslp$
menyatakan diameter tulangan desak pada daerah lapangan, $ntldslp$
menyatakan jumlah tulangan desak pada daerah lapangan, $diatlgs$ menyatakan
diameter tulangan geser, $jaraktlgs$ menyatakan jarak antar tulangan geser.
$M$ adalah momen dan $G$ adalah gaya geser.

\medskip\noindent
\textbf{untuk kolom}, minimumkan:
\begingroup\footnotesize
\begin{equation}
\begin{aligned}
f_2(B,H,dia,N,\mathit{diatlgskl},\mathit{jaraktlgskl}) =\;
  &(B \times H \times \mathit{tinggi\ kolom} \times \mathrm{Rp}(\text{harga beton})/\mathrm{m}^3)\\
 +&\left(\tfrac14 \pi\, dia^2 \times (4N-4) \times \mathit{tinggi\ kolom} \times (\text{berat tulangan})_{\mathrm{kg/m}^3} \times \mathrm{Rp}(\text{harga tulangan})/\mathrm{kg}\right)\\
 +&\Big(\tfrac14 \pi\, \mathit{diatlgskl}^2 \times \big(2(B - \mathit{selimut\ beton}) + 2(H - \mathit{selimut\ beton})\big)\\
   &\times (\mathit{tinggi\ kolom}/\mathit{jaraktlgskl}) \times (\text{berat tulangan})_{\mathrm{kg/m}^3} \times \mathrm{Rp}(\text{harga tulangan})/\mathrm{kg}\Big)
\end{aligned}
\label{eq:1-8}
\end{equation}
\endgroup
yang memenuhi kendala persamaan:
\begin{align}
M_{\text{intern}} &\geq M_{\text{extern}} \label{eq:1-9}\\
G_{\text{intern}} &\geq G_{\text{extern}} \label{eq:1-10}\\
P_{\text{intern}} &\geq P_{\text{extern}} \label{eq:1-11}
\end{align}
pada persamaan di atas $f_2$ menyatakan harga kolom pada struktur, dengan
variabel desain $B$ dan $H$ menyatakan ukuran penampang beton kolom, $dia$
menyatakan diameter tulangan, $N$ merupakan jumlah tulangan yang digunakan,
$diatlgskl$ merupakan diameter tulangan geser kolom, $jaraktlgskl$
menyatakan jarak antara tulangan geser pada kolom. $M$ adalah momen, $G$
adalah gaya geser, $P$ adalah gaya aksial pada kolom.

Harga struktur portal seluruhnya dapat dihitung dengan
\begin{equation}
f = \sum_{b=1}^{nB} f_{1b} + \sum_{k=1}^{nK} f_{2k}
\label{eq:1-12}
\end{equation}
pada persamaan di atas $f$ merupakan harga portal total, $f_{1b}$
menyatakan harga balok $b$, $f_{2k}$ menyatakan harga kolom $k$, $nB$
menyatakan jumlah balok, $nK$ menyatakan jumlah kolom.

Untuk menyelesaikan masalah tersebut akan dibuat program komputer yang
ditulis dalam bahasa pemrograman yang sudah banyak digunakan yaitu bahasa
C++ dan menggunakan data diskrit serta menggunakan metoda optimasi
polihedron fleksibel yang dapat digunakan untuk menangani masalah
non-linear.

\subsection{Batasan Masalah}

Dalam melakukan optimasi beton bertulang pada struktur portal ruang ini
perhitungan beton bertulang didasarkan pada peraturan yang berlaku di
Indonesia yaitu Tata Cara Perhitungan Struktur Beton Untuk Bangunan Gedung
SK SNI T-15-1991-03 (PU, 1991), sedangkan untuk melakukan perhitungan gaya
pada struktur digunakan metoda kekakuan (Weaver, 1980). Mengingat luasnya
permasalahan dan keterbatasan waktu, maka dalam tugas akhir ini diperlukan
pembatasan masalah, yaitu:
\begin{enumerate}[label=\alph*.]
\item Analisa struktur portal ruang memakai metoda kekakuan
\item Rigid-end offset diabaikan
\item Beban yang bekerja pada struktur adalah beban gravitasi menerus
\item Balok dan kolom berpenampang empat persegi panjang dan uniform
      sepanjang bentangnya
\item Tidak ada gelincir antara beton dengan tulangan baja
\item Perencanaan kolom biaksial menggunakan metoda kontur beban
\item Gaya puntir diabaikan
\item Perencanaan beton bertulang menggunakan analisis ultimit
\end{enumerate}

\section{Maksud dan Tujuan Penulisan}

Penulisan tugas akhir ini dimaksudkan untuk memenuhi syarat yudisium dalam
mencapai tingkat kesarjanaan Strata 1 (S1) pada Program Studi Teknik Sipil,
Fakultas Teknik, Universitas Atma Jaya Yogyakarta.

Selain itu penulisan tugas akhir ini juga bertujuan untuk mengetahui lebih
dalam tentang metoda-metoda optimasi yang sekarang ini mulai dirasakan
manfaatnya dan mulai dikembangkan, terutama metoda yang dapat digunakan
untuk melakukan optimasi biaya portal beton bertulang, yang memiliki
permasalahan yang kompleks.

Hasil dari penulisan tugas akhir ini adalah berupa program komputer yang
ditulis dengan bahasa C++ yang dapat dikompilasi menggunakan Borland C++
5.0 atau versi yang lebih tinggi pada sistem operasi Windows 32 bit menjadi
sebuah perangkat lunak.

\section{Kegunaan Penulisan}

Perangkat lunak yang dihasilkan dalam penulisan ini berguna untuk
merencanakan ukuran penampang balok dan kolom beserta penulangannya dengan
memasukkan informasi tentang bentuk geometri portal ruang yang akan
direncanakan.

Selain itu program yang dibuat ini dapat dijadikan dasar pengembangan
perangkat lunak untuk menyelesaikan masalah optimasi sehingga dapat
tercipta sebuah program yang lebih baik lagi, sehingga akan didapatkan
hasil yang lebih optimum lagi.

\section{Metoda Penulisan}

Dalam melakukan penulisan ini banyak langkah-langkah yang diambil sebelum
mendapatkan sebuah perangkat lunak yang siap digunakan untuk melakukan
optimasi.

Tahap pertama, dilakukan studi pustaka mengenai optimasi beton bertulang
pada portal ruang, metoda yang digunakan untuk menganalisis struktur
digunakan metoda kekakuan, beton dianalisis dengan analisis ultimit,
sedangkan kolom biaksial dianalisis menggunakan metoda kontur beban.

Tahap kedua adalah membuat suatu listing program untuk melakukan analisa
struktur menggunakan metoda kekakuan, juga membuat listing program untuk
melakukan analisa balok menggunakan analisis ultimit, serta listing program
untuk melakukan analisis kolom biaksial menggunakan metoda kontur beban.

Tahap ketiga adalah menggabungkan listing program yang telah diuji
validasinya masing-masing sebagai subrutin yang dipanggil oleh sebuah
fungsi utama, dalam fungsi utama itu dilakukan proses optimasi.

Tahap keempat adalah melakukan uji validasi dari perangkat lunak yang
digunakan dengan hasil yang pernah dihitung tanpa menggunakan proses
optimasi atau dengan perangkat lunak lain yang sudah jadi.

Tahap terakhir adalah mencoba perangkat lunak untuk memecahkan kasus-kasus
yang ada untuk melakukan penyusunan laporan tugas akhir secara lengkap.
